% !TeX program = xelatex
% Edit this file directly; cv-style.sty controls the layout.
\documentclass[10pt,a4paper]{article}
\usepackage{cv-style}
% Shared contact details and content review date for both languages.
% Change the date when the content is reviewed, not just when recompiling.
\newcommand{\CVUpdated}{2026-09-06}
\newcommand{\CVEmail}{23300240019@m.fudan.edu.cn}
\newcommand{\CVPhone}{+86 158-0701-0023}
\newcommand{\CVWebsite}{https://16yunh.github.io/}
\newcommand{\CVWebsiteLabel}{16yunh.github.io}
\newcommand{\CVGitHub}{https://github.com/16yunH}
\newcommand{\CVGitHubLabel}{github.com/16yunH}
\newcommand{\CVPortrait}{assets/portrait.jpg}

\hypersetup{pdftitle={Yun Hong - Curriculum Vitae},pdflang={en}}
\begin{document}
\cvsetup{en}
\cvheader{Yun Hong}{COMPUTER SCIENCE \textbar{} FUDAN UNIVERSITY}
I conduct research at MagicLab, Fudan University, under the guidance of Wenchao Ding. My current research focuses on world models for embodied intelligence.

\cvsection{Education}
\begin{cventry}{Fudan University}{Sep 2023 - Present}{B.S. candidate in Computer Science \textbar{} Shanghai, China}
\end{cventry}

\begin{cventry}{The University of Texas at Austin}{Aug - Dec 2025}{Exchange student in Computer Science \textbar{} Austin, TX, USA}
\begin{cvitems}
\item GPA: 4.0/4.0 (13 credits). A grades in machine learning, computer vision, operating systems, and probability/random processes.
\end{cvitems}
\end{cventry}

\cvsection{Research Experience}
\begin{cventry}{World Models for Embodied Intelligence}{Current research}{MagicLab, Fudan University \textbar{} Advisor: Wenchao Ding}
\begin{cvitems}
\item Investigating how execution feedback can help embodied world-action models detect deviations and correct subsequent actions.
\end{cvitems}
\end{cventry}

\begin{cventry}{Risk-Map Learning for Autonomous Driving}{Oct 2024 - May 2025}{MagicLab, Fudan University \textbar{} Advisor: Wenchao Ding}
\begin{cvitems}
\item Contributed to the Transformer-based risk-prediction framework, model training, debugging, and result verification.
\item Fourth author of the RA-L 2026 publication listed below.
\end{cvitems}
\end{cventry}

\begin{cventry}{Multi-Task Active Perception}{Fall 2025}{RobIn Lab, UT Austin \textbar{} Advisor: Junhong Xu}
\begin{cvitems}
\item Built multi-task partially observable simulation benchmarks involving unknown object properties and spatial locations.
\item Developed VLA baseline training and evaluation workflows to study imitation learning, reinforcement learning, and IL pretraining followed by RL fine-tuning.
\item Reviewed manipulation exploration and long-horizon POMDP benchmarks; the project primarily completed benchmark and baseline evaluation workflows.
\end{cvitems}
\end{cventry}

\begin{cventry}{LLM Reasoning, Reflection, and Self-Correction}{Sep 2024 - Jun 2025}{Fudan NLP \textbar{} Advisor: Xuanjing Huang}
\begin{cvitems}
\item Constructed self-critic data using Llama3.1-8B-Instruct, with tree-search and reverse-reasoning strategies for diverse correction trajectories.
\item Built data synthesis and training workflows; evaluated reasoning and reflection methods on GSM8K.
\end{cvitems}
\end{cventry}

\cvsection{Publication}
% Shared publication record; maintain author order here.
{\cvsmall
\textbf{Learning a Unified Risk Map for Autonomous Driving in Partially Observable Environments}\par
Jie Jia, Yaofeng Su, Zeyu Bao, \textbf{Yun Hong}, Bingzhao Gao, Zhongxue Gan, Wenchao Ding\par
\textit{IEEE Robotics and Automation Letters}, 11(4): 4457-4464, 2026.\par
\href{https://doi.org/10.1109/LRA.2026.3666393}{doi:10.1109/LRA.2026.3666393}\par}


\newpage
\cvsection{Internship Experience}
\begin{cventry}{ByteDance}{Jan - Jun 2026}{Test Development Engineer Intern \textbar{} Douyin R\&D, Live Streaming \& Media Stream QA}
\begin{cvitems}
\item Developed three AI-assisted tools for routine monitoring, test-case risk checking, and alert summarization.
\item Refactored and expanded inspection and test cases. Monitoring accuracy increased from 89\% to 100\% and coverage from 95\% to 98\% within the evaluated scope.
\end{cvitems}
\end{cventry}

\cvsection{Selected Projects}
\begin{cventry}{LCDP-Sim: Language-Conditioned Diffusion Policy}{Started Dec 2025}{\href{https://github.com/16yunH/LCDP-Sim}{https://github.com/16yunH/LCDP-Sim}}
\begin{cvitems}
\item Built a vision-language-action system with CLIP and a U-Net DDPM/DDIM diffusion policy, including 16-step action chunking.
\item Implemented data collection, training, and closed-loop evaluation workflows in ManiSkill2.
\end{cvitems}
\end{cventry}

\begin{cventry}{MapNavigation}{Started Dec 2024}{\href{https://github.com/16yunH/MapNavigation}{https://github.com/16yunH/MapNavigation}}
\begin{cvitems}
\item Integrated OpenStreetMap and Gaode Maps API for route planning, location search, hybrid point selection, and road-type speed constraints.
\end{cvitems}
\end{cventry}

\begin{cventry}{NeuralStyle}{Started Jun 2025}{\href{https://github.com/16yunH/NeuralStyle}{https://github.com/16yunH/NeuralStyle}}
\begin{cvitems}
\item Developed a modular Python style-transfer toolkit with batch processing, configurable experiments, and a browser-based interface.
\end{cvitems}
\end{cventry}

\cvsection{Skills}
\cvdetail{Machine learning}{Supervised/unsupervised learning, reinforcement learning, imitation learning, model evaluation and generalization.}
\cvdetail{Deep learning \& robotics}{PyTorch, Transformers, diffusion models, LLaMA fine-tuning, VLA policies, active perception, POMDPs, and robotic manipulation.}
\cvdetail{Vision \& driving}{Multi-view geometry, 3D reconstruction, detection and segmentation, motion planning, trajectory prediction, and risk-field modeling.}
\cvdetail{Programming \& languages}{C, C++, Python, LaTeX; Chinese (native), English (fluent).}
\cvsection{Honors \& Awards}
\cvdetail{2026}{Award recipient, Tencent Rhino-Bird Open-Source Talent Program: OpenCloudOS Security Skill issue (one of three awardees).}
\cvdetail{2025}{Outstanding Work Award, The 5th Meituan Business Analysis Elite Competition.}
\cvdetail{2024}{Shanghai Third Prize, China Undergraduate Mathematical Contest in Modeling (CUMCM).}
\cvdetail{2024}{Finalist, Water, sanitation, and hygiene for the prevention and care of neglected tropical diseases (Geneva, Switzerland).}
\cvdetail{2023}{Finalist, Full-stack AI development engineer skills training by NVIDIA.}
\end{document}
